\appendix

\section*{Appendix}

\section{Retrieval Algorithms}
\label{app:algorithms}

\begin{algorithm}[h]
\caption{Context Retrieval from Knowledge Base}
\label{alg:context-retrieval}
\begin{algorithmic}[1]
\REQUIRE Task description $q$; knowledge base $\mathcal{K}$; neighbor count $k{=}5$; scaling factor $\alpha{=}1.5$; cap $p \in [5, 25]$; validation weights $W$
\ENSURE Selected learnings $C$

\STATE $\mathbf{v}_q \leftarrow \textsc{Embed}(q)$
\STATE $\mathcal{T}_{\text{top}} \leftarrow \text{Top-}k \text{ tasks in } \mathcal{K} \text{ by } \text{sim}(\mathbf{v}_q, \mathbf{v}_t)$ \COMMENT{filter exact matches}
\STATE $L_{\max} \leftarrow \max_{t \in \mathcal{T}_{\text{top}}} \textsc{LearningCount}(t)$
\STATE $N_{\text{pick}} \leftarrow \min(\lceil \alpha \cdot L_{\max} \rceil, \; p)$
\STATE $\mathcal{L}_{\text{pool}} \leftarrow \emptyset$
\FOR{each $(t, s_t) \in \mathcal{T}_{\text{top}}$}
    \FOR{each learning $\ell \in \textsc{GetLearnings}(t)$}
        \STATE $\text{score}(\ell) \leftarrow s_t \times W[\ell.\text{valid\_level}]$
        \STATE $\mathcal{L}_{\text{pool}} \leftarrow \mathcal{L}_{\text{pool}} \cup \{(\ell, \text{score}(\ell))\}$
    \ENDFOR
\ENDFOR
\STATE Sort $\mathcal{L}_{\text{pool}}$ by score descending; deduplicate by learning text
\STATE $C \leftarrow$ first $N_{\text{pick}}$ elements of $\mathcal{L}_{\text{pool}}$
\RETURN $C$
\end{algorithmic}
\end{algorithm}

\begin{algorithm}[h]
\caption{Hybrid Tool Retrieval}
\label{alg:tool-retrieval}
\begin{algorithmic}[1]
\REQUIRE Issue description $I_q$; knowledge base $\mathcal{K}$; result count $k{=}3$; dense weight $\alpha{=}0.7$; sparse weight $\beta{=}0.3$; validation weights $W$
\ENSURE Top-$k$ issue--learning pairs

\STATE $q' \leftarrow \textsc{Clean}(I_q)$ \COMMENT{remove digits, normalize}
\STATE $\mathbf{v}_q \leftarrow \textsc{Embed}(q')$
\STATE $\mathbf{s}_{\cos} \leftarrow \textsc{CosineSim}(\mathbf{v}_q, \mathcal{K}.\text{issue\_embeddings})$
\STATE $\mathbf{s}_{\text{bm25}} \leftarrow \textsc{BM25}(q', \mathcal{K}.\text{index})$
\STATE $\mathcal{C} \leftarrow \text{Top}(\mathbf{s}_{\cos}, 100) \cup \text{Top}(\mathbf{s}_{\text{bm25}}, 100)$ \COMMENT{candidate pool}
\FOR{each $i \in \mathcal{C}$}
    \STATE $c_i \leftarrow \text{clamp}(\mathbf{s}_{\cos}[i], 0, 1)$
    \STATE $b_i \leftarrow \textsc{MinMaxNorm}(\mathbf{s}_{\text{bm25}}[i], \mathcal{C})$
    \STATE $S_{\text{hybrid}} \leftarrow \alpha \cdot c_i + \beta \cdot b_i$
    \STATE $\text{score}(i) \leftarrow S_{\text{hybrid}} \times W[\mathcal{K}[i].\text{valid\_level}]$
\ENDFOR
\STATE Sort $\mathcal{C}$ by score descending
\RETURN first $k$ elements of $\mathcal{C}$
\end{algorithmic}
\end{algorithm}

\section{Knowledge Base Schema}
\label{app:kb-schema}

Each entry in the knowledge base is a structured record containing the following fields:

\begin{itemize}
    \item \textbf{Task metadata}: task description, object type, action verbs, goal phase (\textsc{Search}, \textsc{Acquire}, \textsc{Transform}, \textsc{Place}, or \textsc{Recover}).
    \item \textbf{Issue--learning pair}: \texttt{issue\_text} (the specific mistake encountered) and \texttt{learning\_text} (the successful correction).
    \item \textbf{Validation level}: a confidence label reflecting downstream outcome:
    \begin{itemize}
        \item \textsc{Valid\_Next\_Trial}: the learning was applied and led to success in a subsequent trial (highest confidence). Weight $\omega_{\text{next}} = 1.0$.
        \item \textsc{Valid\_Same\_Trial}: the learning is associated with success within the same trial. Weight $\omega_{\text{same}} = 0.9$.
        \item \textsc{Candidate}: the learning was generated but the task was not completed (lowest confidence). Weight $\omega_{\text{cand}} = 0.6$.
    \end{itemize}
    \item \textbf{Provenance}: references to the source task, trial number, and step range for both the issue and its resolution.
\end{itemize}

\section{Experimental Details}
\label{app:experimental-details}

\subsection{Dataset Statistics}

\paragraph{ALFWorld.}
The full ALFWorld dataset~\citep{shridhar2021alfworld} comprises 3{,}827 games split into a training set (3{,}553 games), a validation-seen set (140 games), and a validation-unseen set (134 games). For knowledge base construction, we use \emph{alfworld-mini}, a curated subset of 120 training games (20 per task type $\times$ 6 task types). We run Reflexion on this subset and extract approximately 400 issue--learning entries from the resulting trajectories using the pipeline described in Section~\ref{sec:kb-construction}. The validation splits are identical to those of the full ALFWorld dataset, ensuring direct comparability with prior work. No training games appear in the validation sets.

\paragraph{WebShop.}
We partition the WebShop environment~\citep{yao2022webshop} into a development split (200 tasks) and a held-out test split (200 tasks). The knowledge base is extracted from Reflexion training trajectories on a separate set of training tasks with no overlap with either evaluation split, yielding approximately 677 issue--learning entries.

\paragraph{SQL diagnostic gate.}
We include a paired SQL diagnostic gate of 50 tasks run with GPT-5.4 to probe transfer behavior in a formally regular action space. This gate is used as diagnostic evidence rather than as a full third benchmark. The paired run contains zero retrieved-help calls across memory variants and no environment-error matches in world logs, isolating memory and prompt effects from runner instability. Summary outputs are stored in the SQL transfer-radius experiment report directory.

\subsection{Framework Variants}

We evaluate six framework variants in each environment:
\begin{enumerate}
    \item \textbf{ReAct}: the base agent with no memory~\citep{yao2023react}, serving as the primary baseline.
    \item \textbf{ReAct + Reflexion}: the iterative self-correction baseline~\citep{shinn2023reflexion}, allowed up to 7 trials per task.
    \item \textbf{ReAct + CR}: base agent augmented with Context Retrieval only.
    \item \textbf{ReAct + TR}: base agent augmented with Tool Retrieval only.
    \item \textbf{ReAct + CR + TR}: base agent augmented with both retrieval mechanisms.
    \item \textbf{ReAct + hard-neg CR + TR}: ablation control that retrieves maximally \emph{irrelevant} entries via inverted retrieval (Section~\ref{sec:hard-neg}).
\end{enumerate}
All memory-augmented variants (CR, TR, CR+TR) execute a single trial per task, whereas Reflexion may execute up to 7 trials.

\subsection{Evaluation Metrics}

\paragraph{Task success.}
The fraction of tasks completed successfully, computed as a binary outcome per task. For WebShop, this is corrected strict success.

\paragraph{Average Steps per Trial.}
The mean number of thought--action--observation steps per individual trial attempt. For Reflexion, this reflects the average length of a single trial, irrespective of retry count.

\paragraph{Average Steps per Task.}
The mean total steps per task across all trials. For single-trial methods, this equals Steps/Trial. For Reflexion, this accumulates steps across all retry trials (up to 7), capturing the full computational cost.

\paragraph{Gain from ReAct.}
The percentage-point difference between a variant's task success and the ReAct baseline task success on the same split.

\paragraph{SQL reward and regressions.}
For the SQL diagnostic, we report exact task success, average environment reward, and paired regressions or improvements relative to ReAct on the same 50 tasks.

\subsection{Hyperparameters}
\label{app:hyperparameters}

Table~\ref{tab:hyperparameters} lists all hyperparameters used across experiments.

\begin{table}[h]
\centering
\small
\begin{tabular}{lr p{2.7cm}}
\toprule
\textbf{Parameter} & \textbf{Value} & \textbf{Description} \\
\midrule
\multicolumn{3}{l}{\emph{Context Retrieval}} \\
$p$ & $[5, 25]$ & Tunable cap on retrieved learnings \\
$\alpha$ & 1.5 & Dynamic context sizing factor \\
$k_{\text{CR}}$ & 5 & Neighbor tasks retrieved \\
\midrule
\multicolumn{3}{l}{\emph{Tool Retrieval}} \\
$k_{\text{TR}}$ & 3 & Results returned per query \\
Dense weight & 0.7 & Cosine similarity fusion weight \\
Sparse weight & 0.3 & BM25 fusion weight \\
Candidate pool & $100{+}100$ & Top-100 from each strategy \\
\midrule
\multicolumn{3}{l}{\emph{Validation Weights}} \\
$\omega_{\text{next}}$ & 1.0 & \textsc{Valid\_}\newline\textsc{Next\_Trial} \\
$\omega_{\text{same}}$ & 0.9 & \textsc{Valid\_}\newline\textsc{Same\_Trial} \\
$\omega_{\text{cand}}$ & 0.6 & \textsc{Candidate} \\
\midrule
\multicolumn{3}{l}{\emph{Agent Configuration}} \\
Max steps/task & 49 & Episode length limit \\
Temperature & 0.0 & Deterministic decoding \\
Max output tokens & 2{,}048 & LLM generation limit \\
Max Reflexion trials & 7 & Retry budget \\
\bottomrule
\end{tabular}
\caption{Hyperparameters used across all experiments.}
\label{tab:hyperparameters}
\end{table}

\subsection{Implementation Details}
\label{app:implementation-details}

All agent reasoning uses Gemini 2.5 Flash routed through the OpenRouter API. We use greedy decoding (temperature 0.0) with a maximum output length of 2{,}048 tokens. API calls employ exponential backoff (2\,s--30\,s) with up to 8 retry attempts.

Retrieval embeddings use \texttt{gemini-embedding-001} via the Google GenAI SDK directly, with a batch size of 100 texts per API call. Embeddings are cached to disk alongside the knowledge base file; the cache includes a hash of the source file and the embedding model name for automatic invalidation on changes.

Each task trial is limited to 49 thought--action--observation steps across all agent variants. For Reflexion, this limit applies per trial, with up to 7 trials per task.
